% ============================================================
% 04_results.tex — Резултати и евалуација (македонски)
% ============================================================

\section{Метрики за евалуација на моделот}

Евалуацијата на системот се спроведува на три нивоа: точност на поединечните ML модели, перформанси на сезонскиот симулатор и квалитативна анализа на логиката на одлучување.

\textbf{Примарна метрика}: Средна апсолутна грешка (MAE) врз предикциите на поени по играч и по натпреварувачка недела. MAE е претпочитана поради нејзината интерпретабилност во доменот на FPL поените и поради нејзината робусност спрема исклучително високи поединечни перформанси.

\textbf{Секундарна метрика}: Точност на врвните-$N$ (Top-$N$ Accuracy), која мери дали моделот ги рангира вистинските врвни-$N$ перформери во рамките на врвните предвидени позиции за дадена позициска група. Оваа метрика е подиректно поврзана со практичниот квалитет на изборот што го прави ILP оптимизаторот.

\textbf{Терциерна метрика}: Резултати за важност на карактеристиките (LGBM split count и gain), кои даваат увид во тоа кои сигнали се најинформативни низ позициите.

\section{Резултати од walk-forward валидацијата}

Walk-forward валидацијата низ петте набори даде конзистентни вредности на MAE. Поновите набори (4 и 5) носат поголем пондер во пондерираната просечна евалуација затоа што се порепрезентативни за целната задача на предвидување (предвидување на перформансите за 2025--26).

MAE значително варира по позиција: голманите покажуваат пониска MAE (помала варијанса во поените — главно условена од clean sheet и одбрани), додека средните играчи и нападачите покажуваат повисока MAE (поттикната од високата варијанса на голови и асистенции). LightGBM конзистентно постигнува пониска MAE од XGBoost низ сите пет набори и сите четири позиции. Подобрувањето е најизразено кај средните играчи и нападачите, каде што способноста на стратегијата за раст по листови (leaf-wise growth) да фати нелинеарни интеракции во просторот на карактеристиките е најкорисна.

% Placeholder за графикон на важноста на карактеристиките
\begin{figure}[h]
\centering
\fbox{\parbox{12cm}{\centering
Важност на карактеристиките на LightGBM — MID модел\\
\small (Placeholder — вметни LGBM feature importance графикон тука)\\
X-оска: вредност на важноста (gain); Y-оска: име на карактеристика
}}
\caption{Важност на карактеристиките на LightGBM моделот за средните играчи (MID). Карактеристиките за тековна форма (\texttt{form\_last3}, \texttt{form\_last5}) и тимскиот xG (\texttt{team\_xG\_last5}) се рангираат меѓу најинформативните предиктори.}
\label{fig:feature_importance_mk}
\end{figure}

\section{Перформанси на сезонскиот симулатор}

\subsection{Табела со резултати за GW1--28}

Финалниот систем (Optuna Trial 429) беше евалуиран низ GW1--28 од сезоната 2025--26 на FPL. Табела~\ref{tab:gw_results_mk} ја прикажува целосната распределба по натпреварувачка недела.

\begin{table}[h]
\centering
\caption{Резултати по натпреварувачка недела (GW1--28), сезона 2025--26}
\label{tab:gw_results_mk}
\begin{tabular}{clp{6.5cm}}
\toprule
\textbf{ГН} & \textbf{Поени} & \textbf{Капитен / Чип} \\
\midrule
1  & 79  & Haaland (C) \\
2  & 83  & Haaland (C) \\
3  & 63  & Junqueira de Jesus (C) \\
4  & \underline{98}  & Haaland (C) \\
5  & 62  & Haaland (C) \\
6  & \textbf{101} & Haaland (C) — Triple Captain \\
7  & 70  & Haaland (C) \\
8  & 91  & Haaland (C) — Bench Boost \\
9  & 47  & Haaland (C) \\
10 & 74  & Haaland (C) \\
11 & 46  & Haaland (C) \\
12 & 45  & Lomba Neto (C) \\
13 & 61  & Haaland (C) \\
14 & 60  & Borges Fernandes (C) \\
15 & 61  & Haaland (C) \\
16 & 82  & Borges Fernandes (C) \\
17 & 65  & Borges Fernandes (C) — Wildcard \\
18 & 38  & Haaland (C) \\
19 & 50  & Semenyo (C) \\
20 & 38  & Garner (C) \\
21 & 100 & Nascimento Rodrigues (C) — Triple Captain \\
22 & 42  & Garner (C) \\
23 & 63  & Nascimento Rodrigues (C) — Bench Boost \\
24 & 56  & Borges Fernandes (C) \\
25 & 62  & Borges Fernandes (C) \\
26 & 60  & Borges Fernandes (C) — Free Hit \\
27 & 47  & Junqueira de Jesus (C) \\
28 & 55  & Palmer (C) \\
\midrule
\textbf{Вкупно} & \textbf{1{,}799} & Просек: 64.3 поени/ГН \\
\bottomrule
\end{tabular}
\end{table}

Највисокиот резултат во една натпреварувачка недела беше \textbf{ГН6 со 101 поен}, постигнат при првото активирање на Triple Captain врз Erling Haaland. Вториот по висина беше \underline{ГН4 со 98 поени}. Најниските резултати беа во ГН18 и ГН20 (по 38 поени), и двата во тежок период од средината на сезоната со предизвикувачки натпревари и пад на формата. Системот собра вкупно \textbf{1{,}799 поени} низ 28 натпреварувачки недели со просек од \textbf{64.3 поени по натпреварувачка недела} и нула казнени поени од трансфери.

\subsection{Контекст во рамките на глобалната FPL заедница}

За да се согледа значењето на 1{,}799 поени, неопходно е резултатот да се спореди со пошироката конкуренција во FPL. Играта привлекува \textbf{над 10 милиони учесници} секоја сезона. Просечниот FPL менаџер постигнува приближно \textbf{50 поени по натпреварувачка недела}, што соодветствува на околу \textbf{1{,}400 кумулативни поени низ 28 натпреварувачки недели} — системот ја надминува таа основа за повеќе од \textbf{400 поени (+28.5\%)}.

Од клучно значење е дека индустриските анализи утврдуваат оти за пласман во \textbf{топ 1{,}000 глобално} — врвните 0.01\% од сите учесници — е потребно одржување на просек од приближно \textbf{64 поени по натпреварувачка недела} низ целата сезона, што е еквивалентно на околу 2{,}450 вкупни поени низ 38 натпреварувачки недели. Системот постигна \textbf{64.3 поени по натпреварувачка недела} низ ГН1--28, што го позиционира токму на — и незначително над — прагот потребен за \textbf{глобален пласман во топ 1{,}000} меѓу над 10 милиони менаџери. Со ова системот се сместува во \textbf{топ 0.01\%} од целата FPL заедница.

Овој резултат е постигнат без никаква рачна интервенција: секоја одлука за трансфер, капитенство и чип се донесува автономно само врз база на јавно достапни податоци (FPL API, Understat, Fantasy Football Scout). Дополнителен показател за стратегискиот квалитет е целосното отсуство на казнени поени од трансфери низ сите 28 натпреварувачки недели — субоптималното водење на трансферите е чест облик на пропаст и кај човечките менаџери и кај наивните автоматизирани пристапи, но беше целосно избегнат преку комбинацијата од логиката за банкирање слободни трансфери, правилата за заклучување на чиповите и бонусите за лојалност во ILP оптимизаторот.

\section{Историја на подобрувањата}

Системот беше развиван итеративно, при што секоја архитектонска компонента беше додавана постепено. Табела~\ref{tab:improvement_mk} го следи развојот од почетната основа до финалниот систем.

\begin{table}[h]
\centering
\caption{Историја на подобрувањата на системот (на размер ГН1--28 освен ако не е поинаку наведено)}
\label{tab:improvement_mk}
\begin{tabular}{lc}
\toprule
\textbf{Конфигурација} & \textbf{Поени} \\
\midrule
Stage 8 основа (без intel, без GW снимки)        & $\approx$ 429 \\
+ GW-snapshot карактеристики                      & $\approx$ 557 \\
+ Intel казни + FDR + бонуси за лојалност         & $\approx$ 616 \\
+ Бонус за сопственост + заклучување на чипови (пред исправка на бубачката) & $\approx$ 654 \\
+ Поправки на TC/BB активирање + исправки на бубачки & 652 (исправено) \\
Целосна сезона XGBoost Trial 1 (ГН1--28)          & 1{,}716 \\
Премин на LGBM Trial 220                          & 1{,}760 \\
Optuna Trial 429 (тековен најдобар)               & \textbf{1{,}799} \\
\bottomrule
\end{tabular}
\end{table}

Преминот од XGBoost кон LightGBM (Trial 1 кон Trial 220) донесе добивка од +44 поени. Последователната Бајесова оптимизација (Trial 220 кон Trial 429) додаде уште +39 поени, што ја покажува вредноста и на изборот на алгоритам и на внимателното подесување на хиперпараметрите.

\section{Резултати од пребарувањето на хиперпараметри}

Заедничкото случајно пребарување со 250 обиди покажа јасна доминација на LightGBM над XGBoost: \textbf{14 од 15-те најдобри конфигурации од обидите користат LightGBM}. Единствената XGBoost конфигурација во топ 15 (Trial 1, основната конфигурација) постигна 1{,}716 поени, додека најдобриот LightGBM обид во оваа фаза (Trial 220) постигна 1{,}760 поени — разлика од 44 поени. Овој наод е конзистентен низ сите модели по позиција и низ сите подмножества на набори.

Фазата на Бајесова оптимизација (100 Optuna обиди, TPE сампл-ер) постигна дополнително подобрување до 1{,}799 поени во Trial 429. Способноста на TPE сампл-ерот да ја моделира распределбата на конфигурациите со добри перформанси му овозможи да конвергира кон оптимумот побрзо од случајното пребарување, во согласност со теоретските предвидувања на Snoek и соработниците \cite{snoek2012practical} и емпириските наоди на Akiba и соработниците \cite{akiba2019optuna}.

\section{Влијание на intel pipeline}

Кумулативното влијание на компонентите на intel pipeline е видливо во историјата на подобрувањата. Додавањето на intel казни за достапност и за ризик од ротација (intel\_03, intel\_04, intel\_06) придонесе со приближно 59 поени над основата со GW-snapshot карактеристики. Оваа добивка ја одразува директната корист од избегнување на играчите кои биле повредени, сомнителни или интензивно ротирани во неделите кога биле избирани.

Бонусот за сопственост во ГН1 (\texttt{OWN\_BOOST\_GW1 = 0.213}) го одразува емпирискиот наод дека на почетокот на сезоната, пред да биде достапен било каков податок за форма, играчите со висока сопственост систематски се потценети од чистите модели за предикција. Тоа е затоа што популарниот консензус често вклучува информации од предсезонските тренинзи и менаџерските намери кои не се фатени во историската статистика.

\section{Анализа на активирањето на чиповите}

\begin{table}[h]
\centering
\caption{Активирања на чипови во текот на сезоната}
\label{tab:chips_mk}
\begin{tabular}{llp{6.5cm}}
\toprule
\textbf{Чип} & \textbf{ГН} & \textbf{Причина за активирање и резултат} \\
\midrule
Triple Captain 1 & ГН6  & Формата на Haaland го надмина TC\_THRESH; постигнати 101 поени \\
Bench Boost 1    & ГН8  & Просечната предикција на клупата го надмина BB\_THRESH; 91 поени \\
Wildcard 1       & ГН17 & Повеќе од 5 членови на составот под просекот за позицијата \\
Triple Captain 2 & ГН21 & Nascimento Rodrigues во одлична форма; 100 поени \\
Bench Boost 2    & ГН23 & Lookahead за клупата активиран од intel\_07; 63 поени \\
Free Hit 1       & ГН26 & Двојна натпреварувачка недела го активира FH; 60 поени \\
\bottomrule
\end{tabular}
\end{table}

Двете активирања на Triple Captain (ГН6: 101 поени, ГН21: 100 поени) беа одлуките за чипови со највисоко влијание. Wildcard во ГН17 беше клучен за ресетирање на составот при влезот во втората половина од сезоната. Активирањата на Bench Boost беа поддржани од оценувањето на кандидатите за клупата воведено во intel\_07, кое обезбеди квалитетот на клупата да биде изграден однапред.

\section{Откриени и поправени бубачки}

Во текот на развојот на системот беа идентификувани и поправени пет критични бубачки:

\begin{table}[h]
\centering
\caption{Критични бубачки идентификувани и поправени во текот на развојот}
\label{tab:bugs_mk}
\begin{tabular}{p{3cm}p{3.5cm}p{3.5cm}c}
\toprule
\textbf{Бубачка} & \textbf{Основна причина} & \textbf{Применета поправка} & \textbf{Поврат} \\
\midrule
Погрешен знак на казната & Казните за трансфери се \textit{собираа} наместо да се одземаат & Одземање со негативен знак & Корекција на резултатот \\
Bench Boost никогаш не се активира & Погрешен клуч во речникот за пребарување на предикцијата за клупата & Рачна пресметка на клупата од состав минус XI & Активирање на чипот \\
Triple Captain никогаш не се активира & Проверката за TC користеше сирова предикција без позицискиот множител & Применување на множителот 1.25x за FWD во проверката за TC & Активирање на чипот \\
Двојно броење на DGW & Вториот ред од DGW тивко го препишуваше првиот & Адитивно акумулирање на двата реда & $\approx$+10 поени \\
Кодирање на имињата на играчите & Имињата со акценти даваа пребарувања со 0 поени & Премин на пребарување по player ID & Неколку поени \\
\bottomrule
\end{tabular}
\end{table}

Овие бубачки ја илустрираат важноста на систематското интеграциско тестирање кај комплексните AI системи со повеќе компоненти. Неколку од бубачките произведуваа резултати кои изгледаат веродостојно но се неточни (на пример, чипови што не се активираат, додека системот навидум функционира нормално), што ги правеше тешки за откривање без насочени unit тестови.

\section{Познати ограничувања}

Системот има четири признаени ограничувања:

\begin{enumerate}
  \item \textbf{Празнина во прес-конференциите за Newcastle}: Скрапер-от intel\_02 ги обработува само клубовите спомнати во насловите на написите на Fantasy Football Scout. Newcastle United (а повремено и други клубови) можат да не се појават како заглавија на секции, со што нивните повредени играчи остануваат неоткриени.

  \item \textbf{Free Hit на празни натпреварувачки недели}: Чипот FH се активира само на двојни натпреварувачки недели. Сценаријата со празни натпреварувачки недели (на пример отказите поради AFCON) не активираат автоматски Free Hit.

  \item \textbf{Шема на продажба и повторно купување}: ILP нема меморија за трансферите од минатата недела, што му дозволува да продаде и повторно да го купи истиот играч во последователни натпреварувачки недели. Беше тестирана казна за повторно купување, но таа го намали вкупниот резултат, па беше оставена како познато ограничување.

  \item \textbf{Отсуство на карактеристики како прокси за замор}: Карактеристиките \texttt{minutes\_last3} и \texttt{minutes\_last5} ги нема во тековниот модел. Нивното вклучување би овозможило подобро моделирање на заморот на играчите за време на густи распореди на натпревари.
\end{enumerate}
